\documentclass[10pt]{article}
\usepackage[T1]{fontenc}
\usepackage[utf8]{inputenc}
\usepackage[french]{babel}
\usepackage{amsmath,amssymb,mathrsfs,mathpazo}
\usepackage{geometry}
\geometry{a4paper,margin=2cm}
\usepackage{graphicx,array,tabularx,tabularray,colortbl,multirow,multicol,booktabs,enumitem,pifont,ulem,fancybox,diagbox,cancel}
\usepackage{pstricks,pstricks-add,pst-plot,pst-tree,pst-node,pst-text,pst-func,pst-math,pst-bspline,pst-3dplot,pst-eucl,pst-solides3d}
\let\pscalculate\undefined
\usepackage{pst-bezier}
\usepackage{tikz,pgfplots,tkz-tab}
\usetikzlibrary{arrows,arrows.meta,patterns,calc,positioning,trees,decorations.markings,shapes.geometric,shapes.misc,shapes.arrows,intersections,decorations.pathmorphing,decorations.pathreplacing,backgrounds,shadows,babel}
\pgfplotsset{compat=1.18}
\usepgfplotslibrary{fillbetween}
\usepackage[np]{numprint}
\usepackage{listings,makecell,pas-tableur,twoopt}
\let\barre\undefined
\usepackage[pstricks]{tablvar}
\definecolor{darkgreen}{RGB}{0,100,0}
\providecommand{\red}{\color{red}}
\providecommand{\blue}{\color{blue}}
\providecommand{\green}{\color{green}}
\providecommand{\black}{\color{black}}
\providecommand{\white}{\color{white}}
\providecommand{\gray}{\color{gray}}
\providecommand{\cyan}{\color{cyan}}
\providecommand{\magenta}{\color{magenta}}
\providecommand{\yellow}{\color{yellow}}
\providecommand{\violet}{\color{violet}}
\providecommand{\orange}{\color{orange}}
\providecommand{\euro}{\texteuro}
\providecommand{\EUR}{\texteuro}
\providecommand{\ieme}{\textsuperscript{e}}
\providecommand{\degres}{\ensuremath{^\circ}}
\providecommand{\degree}{\ensuremath{^\circ}}
\providecommand{\vv}[1]{\overrightarrow{#1}}
\providecommand{\Coord}[1]{#1}
\providecommand{\decosix}{}
\providecommand{\starredbullet}{\ensuremath{\bullet}}
\pagestyle{empty}
\setlength{\parindent}{0pt}
\usepackage{adjustbox}\DeclareUnicodeCharacter{2236}{:}\DeclareUnicodeCharacter{2190}{\ensuremath{\leftarrow}}
\providecommand{\R}{\mathbb{R}}
\providecommand{\N}{\mathbb{N}}
\providecommand{\Z}{\mathbb{Z}}
\providecommand{\Q}{\mathbb{Q}}
\providecommand{\D}{\mathbb{D}}
\providecommand{\C}{\mathbb{C}}
\providecommand{\e}{{\mathrm{e}}}
\providecommand{\vect}[1]{\overrightarrow{#1}}
\providecommand{\vectt}[1]{\overrightarrow{\mathrm{#1}}}
\providecommand{\vvt}[1]{\overrightarrow{\mathrm{#1}}}
\providecommand{\np}[2][]{{#2~#1}}
\providecommand{\vv}[1]{{\overrightarrow{#1}}}
\renewcommand{\d}{{\mathrm{d}}}
\providecommand{\degres}{{^\circ}}
\providecommand{\degre}{{^\circ}}
\providecommand{\blue}{\color{blue}}
\providecommand{\red}{\color{red}}
\providecommand{\enskip}{\quad}
\providecommand{\llbracket}{\lbrack\!\lbrack}
\providecommand{\rrbracket}{\rbrack\!\rbrack}
\providecommand{\euro}{{\text{€}}}
\providecommand{\no}{{\text{n°}}}
\providecommand{\interval}[2]{{[#1;#2]}}
\providecommand{\intervFO}[2]{{[#1;#2[}}
\providecommand{\dotfill}{\ldots}

\providecommand{\ds}{\displaystyle}
\providecommand{\textding}[1]{#1}
\providecommand{\Oij}{\(\mathrm{O};\vec{\imath},\vec{\jmath}\)}
\providecommand{\Oijk}{\(\mathrm{O};\vec{\imath},\vec{\jmath},\vec{k}\)}
\providecommand{\Ouv}{\(\mathrm{O};\vec{u},\vec{v}\)}
\providecommand{\intervFF}[2]{\ensuremath{[#1;#2]}}
\renewcommand{\euro}{\ifmmode\text{\texteuro}\else\texteuro\fi}
\DeclareUnicodeCharacter{00B0}{\ifmmode^{\circ}\else\textdegree\fi}
\DeclareUnicodeCharacter{20AC}{\ifmmode\text{\texteuro}\else\texteuro\fi}
% Énoncés : Éducation nationale. Transcription et corrigés : APMEP, Denis Vergès et contributeurs. Découpage et classement : CoopMaths, Sébastien Lozano, Éric Elter et contributeurs. Adaptation de la présentation : Maths Explorer.
% CC BY-SA : https://coopmaths.fr/www/about/
% https://coopmaths.fr/alea/static/bac/2023/tex/bac_2023_03_sujet1_reunion_1.tex
\begin{document}
\noindent\textit{APMEP / CoopMaths. CC BY-SA. Présentation adaptée par Maths Explorer.}\par\medskip
\medskip

Une entreprise appelle des personnes par téléphone pour leur vendre un produit. 

\begin{itemize}
\item L'entreprise appelle chaque personne une première fois:
\begin{itemize}
\item[$\bullet~~$]la probabilité que la personne ne décroche pas est égale à 0,6 ;
\item[$\bullet~~$]si la personne décroche, la probabilité qu'elle achète le produit est égale à 0,3.
\end{itemize}
\item Si la personne n'a pas décroché au premier appel, on procède à un second appel: 
\begin{itemize}
\item[$\bullet~~$]la probabilité que la personne ne décroche pas est égale à $0,3$ ;
\item[$\bullet~~$]si la personne décroche, la probabilité qu'elle achète le produit est égale à 0,2.
\end{itemize}
\item Si une personne ne décroche pas au second appel, on cesse de la contacter.
\end{itemize}

\medskip

On choisit une personne au hasard et on considère les évènements suivants : 

\begin{description}
\item[ ] $D_1$ : \og la personne décroche au premier appel\fg{} ;
\item[ ] $D_2$ : \og la personne décroche au deuxième appel\fg{} ;
\item[ ] $A$ : \og la personne achète le produit \fg.
\end{description}

\begin{center}\emph{Les deux parties peuvent être traitées de manière indépendante}\end{center}

\begin{minipage}{0.48\linewidth}
\textbf{Partie A}

\begin{enumerate}
\item Recopier et compléter l'arbre pondéré ci-contre.
\item En utilisant l'arbre pondéré, montrer que la probabilité de l'évènement $A$ est 

$P(A) = 0,204$.
\item On sait que la personne a acheté le produit.

Quelle est la probabilité qu'elle ait décroché au premier appel ?
\end{enumerate}
\end{minipage}\hfill
\begin{minipage}{0.48\linewidth}
\begin{center}
\pstree[treemode=R,nodesepA=0pt,nodesepB=2.5pt,treesep = 1cm,levelsep=2.cm]{\TR{}}
{\pstree{\TR{$D_1$~}\taput{\ldots}}
	{\TR{$A$}\taput{0,3}
	\TR{$\overline{A}$}\tbput{\ldots}
	}
\pstree{\TR{$\overline{D_1}$~}\tbput{0,6}}
	{\pstree{\TR{$D_{2}$}\taput{\ldots}}
		{\TR{$A$}\taput{\ldots}
		\TR{$\overline{A}$}\tbput{\ldots}
		}
	\TR{$\overline{D_{2}}$}\tbput{\ldots}
		}
}
\end{center}
\end{minipage}

\bigskip

\textbf{Partie B}

\medskip

On rappelle que, pour une personne donnée, la probabilité qu'elle achète le produit est égale à $0,204$.

\medskip

\begin{enumerate}
\item On considère un échantillon aléatoire de $30$ personnes.

On note $X$ la variable aléatoire qui donne le nombre de personnes de l'échantillon qui achètent le produit.
	\begin{enumerate}
		\item On admet que X suit une loi binomiale. Donner, sans justifier, ses paramètres. 
		\item Déterminer la probabilité qu'exactement 6 personnes de l'échantillon achètent
le produit. Arrondir le résultat au millième.
		\item Calculer l'espérance de la variable aléatoire $X$.
	\end{enumerate}
Interpréter le résultat.
\item Soit $n$ un entier naturel non nul.

On considère désormais un échantillon de $n$ personnes. 

Déterminer la plus petite valeur de $n$ telle que la probabilité qu'au moins l'une des personnes de l'échantillon achète le produit soit supérieure ou égale à $0,99$.
\end{enumerate}

\bigskip

\end{document}
